% 148-16(148쪽) — 두 대각선의 길이가 5, 10 이고 두 대각선이 이루는 각이 $\theta$ 인 사각형 ABCD.
% 스캔 화질이 낮아 TikZ 로 다시 그림. 원문과 같이 B 를 왼쪽 아래, C 를 오른쪽 아래, A 를 왼쪽 위, D 를 오른쪽 위에 두고 안쪽을 연하게 칠했다.
% 라벨은 원문에 있는 A·B·C·D·$5$·$10$·$\theta$ 만 둔다 — 사각형의 넓이(답)는 적지 않는다.
\begin{tikzpicture}[line width=0.5pt, scale=0.78, >=stealth]
  \coordinate (B) at (0,0);
  \coordinate (C) at (2.4,0);
  \coordinate (A) at (0.85,1.35);
  \coordinate (D) at (3.9,1.6);
  \coordinate (I) at (1.63,0.669);
  \fill[violet!14] (A) -- (B) -- (C) -- (D) -- cycle;
  \draw[line width=0.55pt] (A) -- (B) -- (C) -- (D) -- cycle;
  \draw[line width=0.55pt] (A) -- (C);
  \draw[line width=0.55pt] (B) -- (D);
  \draw[line width=0.4pt] ([shift=(-41:0.34)]I) arc (-41:22:0.34);
  \node[font=\footnotesize] at ([shift=(-9:0.62)]I) {$\theta$};
  \draw[densely dotted, line width=0.35pt, <->] ([shift=(112:0.24)]B) to[bend left=10]
    node[midway, below=-2.5pt, font=\footnotesize] {$10$} ([shift=(112:0.24)]D);
  \draw[densely dotted, line width=0.35pt, <->] ([shift=(-131:0.20)]A) to[bend right=12]
    node[midway, below=-2.5pt, font=\footnotesize] {$5$} ([shift=(-131:0.20)]C);
  \node[above left=-3.5pt and -3pt, font=\small] at (A) {$\pt{A}$};
  \node[left=1pt, font=\small] at (B) {$\pt{B}$};
  \node[below right=-3pt and -3pt, font=\small] at (C) {$\pt{C}$};
  \node[above right=-3.5pt and -3pt, font=\small] at (D) {$\pt{D}$};
\end{tikzpicture}
